\documentclass[12pt,a4paper]{article}
\usepackage[margin=2.3cm]{geometry}
\usepackage{fontspec}
\setmainfont{Times New Roman}
\usepackage{setspace}
\singlespacing
\usepackage{titlesec}
\titleformat{\section}{\normalfont\large\bfseries}{}{0em}{}
\titleformat{\subsection}{\normalfont\normalsize\bfseries}{}{0em}{}
\titlespacing{\section}{0pt}{10pt}{4pt}
\titlespacing{\subsection}{0pt}{8pt}{2pt}
\usepackage{hyperref}
\hypersetup{hidelinks}
\setlength{\parindent}{0pt}
\setlength{\parskip}{5pt}
\newenvironment{refs}{\begin{list}{}{\setlength{\leftmargin}{1.5em}\setlength{\itemindent}{-1.5em}\setlength{\itemsep}{1pt}}}{\end{list}}

\begin{document}

{\large\bfseries BE451 Degree Project in Finance --- Thesis Synopsis, Fall 2026}

\textbf{Students:} [Full name (student ID)], [Full name (student ID)] \\
\textbf{Preferred tutor:} [Alvin Chen / no preference]

\vspace{4pt}\hrule\vspace{4pt}

\section{Topic 1: Investment and Payout Policy under a Threatened Profit Cap: Evidence from Sweden's Private Welfare Firms}

\subsection{(i) Tentative research question}

\emph{How do firms change their investment and their payouts when the government threatens to cap the return on the capital they keep in the business, and does the change reverse when the threat is withdrawn?}

Firms tend to postpone investment when policy is uncertain, because waiting has option value (Julio and Yook, 2012; Jens, 2017). A threatened profit cap does something more. If the return on capital kept inside the firm may soon be limited by law, that capital is worth less to the owner, so owners have a reason to pay out more while the cap is pending and to stop once it is off the table. The two channels make different predictions: investment should fall under both, but payouts should rise only under the second. The alternative view is that owners never took the proposal seriously, in which case nothing should happen.

Sweden's private welfare sector is a clean place to test this. Every private provider of tax-financed schooling, social care and personal assistance faced a concrete proposal to cap operating profit at the government borrowing rate plus seven percentage points on operating capital: the Reepalu inquiry (8 November 2016), the government bill (19 March 2018), and the Riksdag's rejection of the bill (7 June 2018). Private health-care providers were explicitly left out of the proposal, so they form a natural comparison group with the same public financing. The threat came back with a directive on schools in June 2022, a new bill in June 2026, and the election of 13 September 2026, where profit caps are a campaign issue. We have found no study of how the firms themselves responded. A zero effect is informative: it would tell us how little a loud but legislatively weak threat costs.

\subsection{(ii) Potential data sources and empirical methods}

\textbf{Data.} Serrano (Swedish House of Finance) gives the annual accounts of every Swedish limited company for 2012--2024: investment, dividends, group contributions, cash, debt, employees, entry, exit and group parent. Treated firms are private firms in education (SNI 85) and social care and assistance (SNI 87--88). The placebo group is private health care (SNI 86), which was excluded from the bill. Controls are matched private service firms. We hand-code the ownership type of the largest welfare groups (private equity or not) from group parents and annual reports, and measure payouts at the group level, since private-equity-owned firms pay out through group contributions to holding companies. The four listed welfare firms (Attendo, Humana, AcadeMedia, Ambea) allow a small event study using LSEG or FinBas. Access: Serrano and FinBas are course resources [verified on (date); about (N) treated firms].

\textbf{Method.} A difference-in-differences event study with firm and industry-by-year fixed effects around the three dates in 2016--2018, comparing treated firms with the health-care placebo and with matched service firms. Outcomes: investment over assets, dividends and group contributions over equity, cash, employment growth, exit and ownership change. Mechanism tests: private-equity-owned versus other treated firms, and firms with pre-threat profits above versus below the proposed cap. Robustness: placebo dates before 2016, exclusion of asylum accommodation (SNI 87.9), alternative matching, and the 2022--2026 school episode as a second test. Main threats and how we handle them: the owner-screening rules that took effect on 1 January 2019 (main window 2016--2018; schools separated from care); the rate shock in the second episode (matching and the placebo group); the small listed sample (the event study is illustrative only). Plan: the 2014--2020 analysis on investment and payouts first, then the mechanism tests and the second episode.

\subsection{(iii) Short literature review}

\textbf{Julio and Yook (2012, Journal of Finance)} show across 48 countries that firms cut investment by 4.8 percent on average in election years and catch up afterwards, consistent with waiting until uncertainty is resolved. We use their event-window design on one dated policy threat inside one country, and add payouts and ownership changes as outcomes. \textbf{Jens (2017, Journal of Financial Economics)} shows that the effect is causal using US gubernatorial elections and placebo tests; our placebo is the health-care sector that the bill excluded. \textbf{Kelly, P\'astor, and Veronesi (2016, Journal of Finance)} show that political uncertainty is priced in option markets and that the price depends on how likely the policy change is; we use this to frame the payout prediction and the event study on the listed firms. \textbf{Bonaime, Gulen, and Ion (2018, Journal of Financial Economics)} show that policy uncertainty reduces merger activity; we use it for the ownership-change outcome. Our contribution is to test both the investment channel and the payout channel on the full population of affected firms, with a placebo built into the bill itself, which election-based studies cannot do.

\subsection*{References}
\begin{refs}
\item Bonaime, Alice, Huseyin Gulen, and Mihai Ion, 2018, Does policy uncertainty affect mergers and acquisitions?, \emph{Journal of Financial Economics} 129, 531--558.
\item Jens, Candace E., 2017, Political uncertainty and investment: Causal evidence from U.S. gubernatorial elections, \emph{Journal of Financial Economics} 124, 563--579.
\item Julio, Brandon, and Youngsuk Yook, 2012, Political uncertainty and corporate investment cycles, \emph{Journal of Finance} 67, 45--83.
\item Kelly, Bryan, \v{L}ubo\v{s} P\'astor, and Pietro Veronesi, 2016, The price of political uncertainty: Theory and evidence from the option market, \emph{Journal of Finance} 71, 2417--2480.
\end{refs}

\newpage

\section{Topic 2: Lower Tax, Smaller Tax Shield: How Sweden's 2019 Corporate Tax Reform Changed the Leverage of Private Firms}

\subsection{(i) Tentative research question}

\emph{How did Sweden's 2019 corporate tax reform change the leverage of private firms, and was the response stronger among firms that risked losing part of their interest deductions?}

The reform did two things at once. The corporate tax rate fell from 22 to 21.4 percent in 2019 and to 20.6 percent in 2021, which makes the tax shield on debt less valuable for every firm. At the same time, net interest expense became deductible only up to 30 percent of taxable EBITDA, with a simplification rule that allows SEK 5 million of net interest regardless, shared within a group of associated companies. Standard trade-off theory predicts less debt from both changes. Ivanov, Pettit, and Whited (2026) find the opposite for small US private firms: after tax cuts, cash flows improve, default risk falls and leverage rises. Sanati and Beyhaghi (2025) find that US firms which lost part of their interest deduction in 2017 reduced both debt and investment. Sweden lets us look at both channels in the same reform and in the full population of limited companies. A zero response on both margins would itself be informative: it would say that Swedish private-firm leverage is not tax-driven.

\subsection{(ii) Potential data sources and empirical methods}

\textbf{Data.} Serrano, all Swedish limited companies 2013--2024: debt, interest expense, tax expense, EBITDA, investment, dividends and group parent. Serrano is accounting data, not tax returns, so we cannot observe exactly which firms were limited. We therefore build a \emph{predicted exposure} measure from 2017--2018 accounts: net interest expense relative to 30 percent of accounting EBITDA and relative to the SEK 5 million allowance, aggregated within groups using Serrano's group identifiers. High-exposure firms are those whose pre-reform interest would have exceeded the new limit. The government's review of the rules (SOU 2024:37, pp.~78--82) reports, using tax-return data, that about 7,100 companies had deductions limited in 2019, that half of the limited companies were real-estate firms, and that the number of firms with net interest close to 30 percent of EBITDA rose from about 3,700 a year before the reform to about 6,100 after it. The exposed group is therefore sizeable, concentrated in one sector, and partly adjusting to the rule, which we treat as heterogeneity and as a check on our exposure measure. Access: Serrano is a course resource [verified on (date); about (N) firm-years].

\textbf{Method.} Main test: a difference-in-differences event study of leverage, new borrowing, interest expense over debt and investment on predicted exposure interacted with year indicators, with firm and industry-by-year fixed effects, and matched low-exposure firms of similar size and industry as controls. Second test: the rate cut, using pre-reform tax payments over assets as the dose, so that firms that paid more tax before the reform gained more from the cut. Threats and how we handle them: exposure is predicted, not observed (we show a narrow and a broad definition and check that the bunching at 30 percent documented in SOU 2024:37 also appears in accounting data); firms can shift interest within groups (exposure is measured at group level); the 2022 rate increases (industry-by-year effects and a placebo in 2016--2018); other 2019 rules aimed at multinationals (excluded). Plan: the exposure analysis first, then the rate cut, then real-estate heterogeneity. Fallback: exposure analysis only.

\subsection{(iii) Short literature review}

\textbf{Ivanov, Pettit, and Whited (2026, Review of Economic Studies)} use US state tax changes and private-firm data and find that leverage rises after tax cuts, because lower taxes raise the productivity of capital and lower credit spreads. We test their prediction on the 2019 and 2021 rate steps. \textbf{Sanati and Beyhaghi (2025, Review of Financial Studies)} use the USD 25 million exemption in the 2017 US reform and find that firms which lost the tax benefit of debt cut debt and investment without replacing debt with equity. Our exposure design follows the same logic with predicted exposure rather than a sharp threshold, since Sweden's SEK 5 million rule is an alternative to the EBITDA rule, not a cut-off. \textbf{Heider and Ljungqvist (2015, Journal of Financial Economics)} find that leverage rises by about 0.4 percentage points per percentage point of higher tax, but does not fall when taxes are cut; this asymmetry is one reason a small response to the rate cut would not be surprising, and why the deductibility channel is our main test. \textbf{Brok (2026, Journal of Financial and Quantitative Analysis; 2024, Review of Finance)} shows that tax avoidance and group structure change how debt responds to taxes, which is why we measure exposure at the group level. Prior work on the Swedish reform includes the government review (SOU 2024:37), which documents the limited firms and the bunching at 30 percent with tax data but finds, descriptively and without a control group, that investment in the most affected groups did not fall (pp.~96--97), and a master's thesis at Uppsala University (Jasem and Sado, 2026) that finds modest negative leverage effects among large exposed firms in a firm-level difference-in-differences for 2016--2023. Relative to these, our contribution is to use the full population of private firms rather than large ones, to assign exposure at the group level, to separate the rate-cut channel from the deductibility channel, and to test the investment result with matched controls.

\subsection*{References}
\begin{refs}
\item Brok, Peter, 2024, Debt and taxes: The role of corporate group structure, \emph{Review of Finance} 28, 2051--2082.
\item Brok, Peter, 2026, Debt and taxes: The role of tax avoidance, \emph{Journal of Financial and Quantitative Analysis} 61, 2007--2032.
\item Heider, Florian, and Alexander Ljungqvist, 2015, As certain as debt and taxes: Estimating the tax sensitivity of leverage from state tax changes, \emph{Journal of Financial Economics} 118, 684--712.
\item Ivanov, Ivan T., Luke Pettit, and Toni M. Whited, 2026, Taxes depress corporate borrowing: Evidence from private firms, \emph{Review of Economic Studies} 93, 2715--2749.
\item Jasem, Shahad, and Gabriel Sado, 2026, Leverage under constraint: Firm responses to Sweden's ATAD-based interest limitation reform, Master's thesis, Uppsala University.
\item Sanati, Ali, and Mehdi Beyhaghi, 2025, How does removing the tax benefits of debt affect firms? Evidence from the 2017 U.S. tax reform, \emph{Review of Financial Studies}, forthcoming, doi:10.1093/rfs/hhaf041.
\item SOU 2024:37, F\"orb\"attrade r\"anteavdragsregler f\"or f\"oretag, Statens offentliga utredningar, 31 May 2024.
\end{refs}

\end{document}
